\documentclass[10pt,a4paper]{amsart}
\usepackage[margin=19mm]{geometry}
\usepackage{amsmath,amssymb,mathtools}
\usepackage[scr=rsfs,cal=euler]{mathalfa}
\usepackage{xfrac}
\usepackage[unicode,hidelinks,bookmarks=false]{hyperref}
\newcommand{\ie}{i.e.\ }
\newcommand{\cf}{cf.\ }
\newcommand{\resp}{respectively\ }
\newcommand{\Subsection}{\subsection}
\providecommand{\nobreakdash}{\nobreak\hbox{-}}
\setlength{\emergencystretch}{2em}
\multlinegap0pt
\usepackage{dsfont}
\usepackage{amssymb}
\usepackage[normalem]{ulem}
\usepackage{enumerate}
\usepackage{cancel}
\usepackage{tikz}
\newtheorem{theorem}{Theorem}
\newtheorem{claim}{Claim}
\newcommand{\EE}{ \mathbb{E} }
\newcommand{\N}{\mathbb{N}}
\renewcommand{\P}{\mathbb{P}}
\renewcommand{\emph}[1]{{\sl #1}}
\newcommand{\ep} {\varepsilon }
\newcommand{\la}{\langle}
\newcommand{\ra}{\rangle}
\title{A stochastic Schauder-Tychonoff type theorem and its applications}
\author{Erika Hausenblas, Ankit Kumar$^\ast$, Jonas M. T\"{o}lle}
\date{Source: arXiv:2602.17285v2 (CC BY 4.0)}
\begin{document}
\maketitle

\noindent\emph{Source and licence.} A stochastic Schauder-Tychonoff type theorem and its applications. Version arXiv:2602.17285v2 (CC BY 4.0), distributed by arXiv under the Creative Commons Attribution 4.0 International licence, http://creativecommons.org/licenses/by/4.0/, as recorded in the arXiv licence metadata for this exact version and shown on the abstract page https://arxiv.org/abs/2602.17285v2. Full licence notice: \texttt{licenses/arxiv-2602.17285-CC-BY-4.0.txt}.\par\smallskip

\noindent\textit{Editorial scope.} Counted unit: the article's theorem LCCthrm (source0.tex lines 3103--3128). The article's complete original proof of it is delivered with it (source0.tex lines 3129--3211, one \texttt{proof} environment of 83 lines). No mathematics is written, repaired or re-derived: the statement and the proof are the article's own lines, in the article's own notation, and no retained line is substituted or re-flowed.  Retained uncounted as the source-local context the proof needs: lines 3215--3222.  Nothing was omitted from any retained span, so the package carries no omission record.  No retained line contains a citation, so the wrapper carries no bibliography and no bibliography entry.  No retained line contains a figure environment, an image-inclusion command or drawing code, so no image is delivered and the package declares no asset dependency.  Editorial formatting only, in addition: an AMS \texttt{amsart} preamble that restates the article's own declarations and macros used by the retained lines, namely the environments \texttt{theorem}, \texttt{claim} on separate counters, together with the standard packages those lines need.  The retained environments print in this document as Theorem 1, Claim 1 in order of appearance, whereas the article prints the same items with the numbering of its own full document.  The complete native source of the article as distributed in the arXiv e-print for this exact version is bundled unchanged as \texttt{sources/194949/source0.tex}.
\begin{claim}\label{claim6.3}
	To show
	\begin{align*}
		\P\bigg(\sup_{k\in[0,2^n],k\text{ odd}}|\xi_{k,n}^i|^2> \frac{2^{(n+1)}a_n^2}{C_1}\bigg)
		\simeq
		\sqrt{\frac{C_1}{\pi} }2^{n/2-1}a_n^{-1}e^{-\frac{2^{(n+1)}a_n^2}{C_1}}.
	\end{align*}
\end{claim}

\begin{theorem}[Lévy--Ciesielski construction in infinite-dimension]\label{LCCthrm}
Let	{$\{\xi^{\,i}_{k,n}\}_{0\le k<2^n,n\in \N}^{i\in \mathbb{I}}$} be a sequence of independent standard normal random variables defined on a probability space 
$(\Omega,\mathcal F,\mathbb P)$.  Let $\{\psi_i\}_{i\in \mathbb{I}}$ be the eigenfunctions of the covariance operator $Q$ corresponding to the eigenvalues $\{\lambda_i\}_{i\in \mathbb{I}}$, and the family $\{\psi_i\}_{i\in \mathbb{I}}$  forms a complete orthonormal system in $U$.
	Let $\{h_{k,n}\}$ be the (orthonormal) Haar wavelet system on $L^2(0,1)$ and let
	\[
	S_{k,n}(t):=\int_0^t h_{k,n}(s)\,ds, \quad t\in[0,1]
	\]
	be the associated Schauder functions. Define partial sums
	\[
{	W_n(t):=\sum_{i\in \mathbb{I}}\sum_{m=0}^{n}\ \sum_{(k,m)\in S_{m}}Q^\frac{1}{2}\psi_i\ \xi^{\,i}_{k,m}S_{k,m}(t),\quad t\in[0,1].}
	\]
	where $S_m=\{(k,m): k \text{ odd, } k\leq 2^m\}.$
	Then, the following hold:
	\begin{enumerate}
		\item $W_n\to W$ \emph{uniformly almost surely} on $[0,1]$, hence $W$ has continuous paths and is defined on the same probability space $(\Omega,\mathcal F,\mathbb P)$. In particular,
			\begin{align*}
			\P \Big(\lim_{m\to\infty}\sup_{t\in[0,1]}\big|W_n(t)-W(t)\big|_U>\ep\Big)=0, \ \ \text{ for all }\ \ \ep>0;
		\end{align*}
		\item $W$ is a centred Gaussian process with covariance
		{	\[
	\mathbb E\big[\la W(t),\psi_i\ra \la W(s), \psi_j\ra ]=\lambda_i \delta_{ij}\min\{s,t\},\quad s,t\in[0,1],
		\]where $\delta_{ij}$ is the Kronecker delta.}
		
		Consequently, $W$ is a Wiener process on $[0,1]$.
	\end{enumerate}
\end{theorem}

\begin{proof}
	Let us assume a sequence of Gaussian random variables $\{\xi_{k,n}^i:i\in \mathbb{I},n\in\N, k \text{ odd,}\break k\leq 2^n\}$. 
	Now, define 
	\begin{equation*}
		h_{1,n}(t)=1,
			\end{equation*}and
			\begin{equation*}
	h_{k,n}(t)=
		\left\{\begin{aligned}
			&2^{(n-1)/2}, \quad &&(k-1)2^{-n}<t\leq k2^{-n},\\
			&-2^{(n-1)/2}, \quad  &&k2^{-n}<t\leq (k+1)2^{-n},\\
			&0, &&  \text{ otherwise},
		\end{aligned}
	\right.
	\end{equation*}for $n\in\N$, $0\leq k\leq 2^n$, $k$ odd. 

Let $S_n=\{(k,n): k \text{ odd, } k\leq 2^n\}$, and define $S=\bigcup\limits_{n\in\N}S_n$. Then, $\{h_{k,n}:(k,n)\in S\}$ is a complete orthonormal system in $L^2(0,1)$. Now, we define
	\[
S_{k,n}(t):=\int_0^t h_{k,n}(s)\,ds\quad (t\in[0,1]),
\] then $S_{k,n}$ satisfies the following properties:
\begin{enumerate}
	\item $|S_{k,n}|_\infty=\sup\limits_{t\in[0,1]}|S_{k,n}(t)|=2^{-(n+1)/2}$ (represents a tent-shaped function of height $2^{-(n+1)/2})$;
	\item for each fixed $n$ and $t$, at most one index $k$ has $S_{k,n}(t)\neq 0$, (supports are disjoint on dyadic intervals).
\end{enumerate}
Now, we consider the sequence of partial sums
\[
W_n(t):=\sum_{i\in\mathbb{I}}\sum_{m=0}^{n}\ \sum_{(k,m)\in S_{m}}Q^\frac{1}{2}\psi_i\ \xi_{k,m}^i S_{k,m}(t),\qquad t\in[0,1].
\]
Let us prove (i). For any positive constant $a_n$, 
\begin{align*}\nonumber
&	\P\Big(\sup_{t\in[0,1]} | W_n(t)-W_{n-1}(t) | _U>a_n\Big)
\\&\nonumber
\leq \P\Big(\sup_{t\in[0,1]} | W_n(t)-W_{n-1}(t) | _U^2>a_n^2\Big)
\\&\nonumber
\leq \P\Bigg(\sup_{t\in[0,1]}\Big(\sum_{i\in\mathbb{I}}Q^\frac{1}{2}\psi_i \xi_{k,n}^iS_{k,n}(t),\sum_{i'\in\mathbb{I}}Q^\frac{1}{2}\psi_{i'} \xi_{k,n}^{i'}S_{k,n}(t)\Big)>a_n^2\Bigg)
\\&\nonumber
\leq \P\Bigg(\sup_{t\in[0,1]}\sum_{i\in\mathbb{I}}\lambda_i  | \psi_i | _U^2 |\xi_{k,n}^i|^2|S_{k,n}(t)|^2>a_n^2\Bigg) 
\\&\nonumber
\leq \P\Big(\sum_{i\in\mathbb{I}}\lambda_i 2^{-(n+1)}\sup_{k}|\xi_{k,n}^i|^2>a_n^2\Big)
 \ \Big(\text{using orthonormality of }\psi_i, (i) \text{ for }S_{k,n}\Big)
\\&
 \sim  \P\left(\sup_{k}|\xi_{k,n}^i|^2>\frac{2^{n+1}a_n^2}{\sum\limits_{i\in\mathbb{I}}\lambda_i }\right)   \quad \Big(\text{using Claim \ref{claim6.3}}\Big)
\\&\nonumber
\leq \sqrt{\frac{C_1}{\pi} }2^{n/2-1}a_n^{-1}e^{-\frac{2^{(n+1)}a_n^2}{C_1}}, \qquad \Big(\text{where }C_1=\sum_{i\in\mathbb{I}}\lambda_i \Big).  \label{}
\end{align*}%
%
Let us choose the constant $a_n$ in such a way that 
{\begin{align}\label{FC}
\sqrt{\frac{C_1}{\pi}}	\sum_{n\in\N}  2^{n/2-1}a_n^{-1}e^{-\frac{2^{(n+1)}a_n^2}{C_1}}<\infty, \text{ and }
	\sum_{n\in\N}a_n<\infty.
\end{align}}
%
%
Let us choose $a_n= \big(\mathfrak{c}n2^{-(n+1)}\big)^{\frac{1}{2}}$, where $\mathfrak{c}>2C_1\ln 2$. Then, conditions change to
\begin{align*}
	\sqrt{\frac{C_1}{2\pi\mathfrak{c}}}\sum_{n\in\N}\frac{(2e^{-\mathfrak{c}/C_1})^n}{n^{\frac{1}{2}}}<\infty, \text{ and }\sqrt{\frac{\mathfrak{c}}{2}}\sum_{n\in\N} \frac{n^{\frac{1}{2}}}{2^{n/2}}<\infty.
\end{align*}
%
The condition \eqref{FC}$_1$ and the Borel-Cantelli lemma ensure that the convergence is almost sure. Therefore, 
\begin{align*}
	\sup_{t\in[0,1]}|W_n(t)-W_{n-1}(t)|_U\leq a_n, \text{ for all large enough }n;
\end{align*}the second condition \eqref{FC}$_2$ guarantees that $W_n$ converges uniformly almost surely to a limit function $W$, therefore, $W$ is continuous.

 From the above arguments, we obtained that $W_n$ converges uniformly to some continuous limit  $W$. Now, we need to show that the limit $W$ is a centred Gaussian process with covariance structure $\mathbb{E}\big[\langle W(t),u\rangle \langle W(s),v\rangle\big]=s\langle Q u,v\rangle,$ for $s\leq t$. Note that, $W_n$ is a centred Gaussian process, thus the vector form $\big(W_n(t_1), W_n(t_2),\dots,W_n(t_k)\big)$ is also a centred Gaussian process and converges almost surely to $\big(W(t_1), W(t_2),\dots,W(t_k)\big)$, which also has a zero-mean Gaussian law. Moreover, the limit of covariances of $W_n$ gives the covariance of $W$
 
Let us move to the proof of covariance. We know that 
\begin{align*}
&	\mathbb{E}\Big[\langle W_n(t),u\rangle \langle W_n(s),v\rangle \Big]\\&=\EE \Bigg[ \Big\langle\sum_{i\in\mathbb{I}}\sum_{m=0}^{n}\ \sum_{(k,m)\in S_{m}}\sqrt{\lambda_i }\psi_i \ \xi_{k,m} ^iS_{k,m}(t),u\Big\rangle\\&\quad \times \Big\langle \sum_{i'\in\mathbb{I}}\sum_{m=0}^{n}\ \sum_{(k,m)\in S_{m}}\sqrt{\lambda_{i'} }\psi_{i'} \ \xi_{k,m}^{i'} S_{k,m}(s),v\Big\rangle\Bigg]
\\&=
\sum_{i\in\mathbb{I}}\langle\sqrt{\lambda_i }\psi_i ,u\rangle\langle \sqrt{\lambda_i }\psi_i ,v\rangle\Bigg(
\sum_{m=0}^{n}\ \sum_{(k,m)\in S_{m}} S_{k,m}(s) S_{k,m}(t)\Bigg)
\\&=
\sum_{i\in\mathbb{I}}\langle\sqrt{Q}\psi_i ,u\rangle\langle \sqrt{Q}\psi_i ,v\rangle\Bigg(
\sum_{m=0}^{n}\ \sum_{(k,m)\in S_{m}} S_{k,m}(s) S_{k,m}(t)\Bigg)
\\&=
\langle Qu,v\rangle\Bigg(
\sum_{m=0}^{n}\ \sum_{(k,m)\in S_{m}} S_{k,m}(s) S_{k,m}(t)\Bigg)
\end{align*}where we used the independence of $\xi_{k,n}$. The right-hand side of the above expression converges to 
\begin{align*}
	\langle Qu,v\rangle
 \sum_{(k,m)\in S} S_{k,m}(s) S_{k,m}(t)= \int_0^1 \mathds{1}_{[0,s]} (r)\mathds{1}_{[0,t]}(r)\, d r=\min\{t,s\},
\end{align*}since $\displaystyle S_{k,m}(t)=\int_0^1\mathds{1}_{[0,t]}(r)h_{k,m}(r)\, dr$ is the Fourier coefficient of $h_{k,n}$ in the representation of $\mathds{1}_{[0,t]}$ in terms of complete orthonormal system $\{h_{k,n}:(k,n)\in S\}$.
\end{proof}

\end{document}
